% Options for packages loaded elsewhere
\PassOptionsToPackage{unicode}{hyperref}
\PassOptionsToPackage{hyphens}{url}
\PassOptionsToPackage{dvipsnames,svgnames,x11names}{xcolor}
%
\documentclass[
  10pt,
]{article}
\usepackage{amsmath,amssymb}
\usepackage{iftex}
\ifPDFTeX
  \usepackage[T1]{fontenc}
  \usepackage[utf8]{inputenc}
  \usepackage{textcomp} % provide euro and other symbols
\else % if luatex or xetex
  \usepackage{unicode-math} % this also loads fontspec
  \defaultfontfeatures{Scale=MatchLowercase}
  \defaultfontfeatures[\rmfamily]{Ligatures=TeX,Scale=1}
\fi
\usepackage{lmodern}
\ifPDFTeX\else
  % xetex/luatex font selection
  \setmainfont[]{DejaVu Serif}
  \setsansfont[]{DejaVu Sans}
  \setmonofont[]{DejaVu Sans Mono}
  \setmathfont[]{Latin Modern Math}
\fi
% Use upquote if available, for straight quotes in verbatim environments
\IfFileExists{upquote.sty}{\usepackage{upquote}}{}
\IfFileExists{microtype.sty}{% use microtype if available
  \usepackage[]{microtype}
  \UseMicrotypeSet[protrusion]{basicmath} % disable protrusion for tt fonts
}{}
\makeatletter
\@ifundefined{KOMAClassName}{% if non-KOMA class
  \IfFileExists{parskip.sty}{%
    \usepackage{parskip}
  }{% else
    \setlength{\parindent}{0pt}
    \setlength{\parskip}{6pt plus 2pt minus 1pt}}
}{% if KOMA class
  \KOMAoptions{parskip=half}}
\makeatother
\usepackage{xcolor}
\usepackage[margin=2.0cm]{geometry}
\usepackage{longtable,booktabs,array}
\usepackage{calc} % for calculating minipage widths
% Correct order of tables after \paragraph or \subparagraph
\usepackage{etoolbox}
\makeatletter
\patchcmd\longtable{\par}{\if@noskipsec\mbox{}\fi\par}{}{}
\makeatother
% Allow footnotes in longtable head/foot
\IfFileExists{footnotehyper.sty}{\usepackage{footnotehyper}}{\usepackage{footnote}}
\makesavenoteenv{longtable}
\setlength{\emergencystretch}{3em} % prevent overfull lines
\providecommand{\tightlist}{%
  \setlength{\itemsep}{0pt}\setlength{\parskip}{0pt}}
\setcounter{secnumdepth}{-\maxdimen} % remove section numbering
% doc_header.tex -- LaTeX preamble additions used by build_docs.sh for every
% Markdown -> LaTeX -> PDF conversion (README and provenance documents of Revision_SuperLoop).
\usepackage{etoolbox}
% longtable may not be loaded by pandoc when a document has no table
\usepackage{longtable}
% bold vectors: \bm -> \symbf under unicode-math (XeLaTeX), bm package otherwise
\makeatletter
\@ifpackageloaded{unicode-math}{\providecommand{\bm}[1]{\symbf{#1}}}{\usepackage{bm}}
\makeatother
% smaller type in long tables so that wide audit tables fit the text width;
% inside a table a line break is allowed after an underscore of plain text
% (pandoc writes it as \_), so that long identifiers wrap inside their column
% rows are set a little apart, so that symbols with indices do not touch
\AtBeginEnvironment{longtable}{\footnotesize\renewcommand{\arraystretch}{1.45}%
  \let\sptunderscore\_\def\_{\sptunderscore\allowbreak}}
\setlength{\LTcapwidth}{\textwidth}
% allow line breaks after slashes and hyphens inside table cells
\usepackage{xurl}
% code blocks: break over-long lines instead of running into the margin.
% Plain ``` blocks are typeset by pandoc as {verbatim}, highlighted ones
% (```bash) as {Highlighting}; both are re-declared as breakable fancyvrb
% environments.  The continuation line is marked with a small hooked arrow.
\usepackage{fvextra}
\RecustomVerbatimEnvironment{verbatim}{Verbatim}{breaklines,breakanywhere,fontsize=\small}
\makeatletter
\@ifundefined{Highlighting}{}{%
  \RecustomVerbatimEnvironment{Highlighting}{Verbatim}{breaklines,breakanywhere,fontsize=\small,commandchars=\\\{\}}}
\makeatother
% inline code (pandoc writes it as \texttt{...}): never hyphenate -- an inserted
% hyphen could be read as part of a file name -- but allow a line break after
% an underscore, slash, comma, colon or equals sign, so that long paths and
% column lists wrap inside the text width and inside table columns.
% Inside a formula (e.g. the name of a notebook symbol in an equation) the same
% command must first leave math mode, which \text of amsmath does.
\ExplSyntaxOn
\tl_new:N \l__spt_code_tl
\cs_new_protected:Npn \__spt_texttt:n #1
  {
    \tl_set:Nn \l__spt_code_tl {#1}
    \regex_replace_all:nnN { ([\/\,\:\=]) } { \1 \c{allowbreak} } \l__spt_code_tl
    \regex_replace_all:nnN { \c{_} } { \c{_} \c{allowbreak} } \l__spt_code_tl
    \group_begin: \ttfamily \hyphenchar \font = -1 \scan_stop: \l__spt_code_tl \group_end:
  }
\RenewDocumentCommand \texttt { m }
  {
    \mode_if_math:TF { \text { \__spt_texttt:n {#1} } } { \__spt_texttt:n {#1} }
  }
\ExplSyntaxOff
% A figure is written in the Markdown source as an image paragraph followed by
% a caption paragraph ("**Figure S1.** ..."), so that the caption is also
% visible on GitHub.  An image that starts a paragraph is centered.
\usepackage{graphicx}% pandoc loads it only for documents that contain an image
\NewCommandCopy\sptincludegraphics\includegraphics
\RenewDocumentCommand\includegraphics{O{} m}{%
  \ifvmode\begin{center}\sptincludegraphics[#1]{#2}\end{center}%
  \else\sptincludegraphics[#1]{#2}\fi}
% keep figures where they are referenced
\usepackage{float}
\floatplacement{figure}{H}
% running header with the document title
\usepackage{fancyhdr}
\pagestyle{fancy}
\fancyhf{}
\fancyhead[L]{\small PROJECT spacetimes, Revision\_SuperLoop}
\fancyhead[R]{\small\thepage}
\renewcommand{\headrulewidth}{0.4pt}
% 16x16 matrices: tighter columns so that a full matrix fits the text width
\setlength{\arraycolsep}{3.2pt}
% The documents of this project contain long tagged symbols, file paths and
% identifiers that cannot be hyphenated.  A paragraph is therefore allowed to
% have looser lines before TeX gives up and lets a line stick out of the margin.
\AtBeginDocument{\tolerance=2000 \emergencystretch=5em \hbadness=2000}
\ifLuaTeX
  \usepackage{selnolig}  % disable illegal ligatures
\fi
\IfFileExists{bookmark.sty}{\usepackage{bookmark}}{\usepackage{hyperref}}
\IfFileExists{xurl.sty}{\usepackage{xurl}}{} % add URL line breaks if available
\urlstyle{same}
\hypersetup{
  pdftitle={The symbols of Revision\_SuperLoop: registry, collisions resolved, and checks},
  pdfauthor={Prepared for Patrick Lee Nash, Ph.D.},
  colorlinks=true,
  linkcolor={blue},
  filecolor={Maroon},
  citecolor={blue},
  urlcolor={blue},
  pdfcreator={LaTeX via pandoc}}

\title{The symbols of Revision\_SuperLoop: registry, collisions
resolved, and checks}
\usepackage{etoolbox}
\makeatletter
\providecommand{\subtitle}[1]{% add subtitle to \maketitle
  \apptocmd{\@title}{\par {\large #1 \par}}{}{}
}
\makeatother
\subtitle{Provenance document \texttt{P1} of PROJECT spacetimes,
Revision\_SuperLoop}
\author{Prepared for Patrick Lee Nash, Ph.D.}
\date{November 16, 2026}

\begin{document}
\maketitle

\hypertarget{what-this-document-is}{%
\section{1. What this document is}\label{what-this-document-is}}

Hard Rule 5 of the author requires that every symbol of PROJECT
spacetimes denote one object, in the whole project, and that the
resolution of every collision be proved by a master provenance of the
symbols, whose symbols are then used everywhere. This document is the
symbol file of the folder \texttt{Revision\_SuperLoop}: it explains how
the registry \texttt{provenance/symbols\_ll.csv} was built (Section 2),
lists the collisions that were found and how each was resolved (Section
3), names the symbols that this folder shares with the folder
\texttt{Revision\_NoGo} (Section 4), states the meaning of the
plain-text names that the programs, the figures and the data files use
(Section 5), reproduces the eight machine checks of the symbols with
their output (Section 6), and contains the complete table of the 87
registered symbols (Section 7). Every symbol that occurs in this
document is defined in Section 7.

\hypertarget{how-the-registry-is-built}{%
\section{2. How the registry is built}\label{how-the-registry-is-built}}

The script \texttt{code/scripts/write\_symbol\_registry.py} writes the
registry, in eight groups: indices, conventions and number sets; the
bi-octonions and the super-loop; the operators and their eigenvalues;
the covariant families and their classification; the manifolds; the
spin-zero superfluid; power counting; groups, Lie algebras, operations
and labels. It has two kinds of entries.

\begin{itemize}
\tightlist
\item
  \textbf{Taken} (2 entries): symbols of the two master registries of
  the project --- \texttt{provenance/symbols.csv} of the main project
  and
  \texttt{Revision\_Deflation\_CJT\_and\_Kohn-Sham/provenance/symbols\_master.csv}
  --- used here with exactly their meaning. The script copies the symbol
  text and the meaning from the master file named in the entry, so that
  they cannot drift, and appends a note on the use here. The taken
  symbols are the index of the space-like coordinates (registered in the
  main project as \(i,j\); here only \(j\) is an index, with the values
  1, 2, 3, and the letter \(i\) is the imaginary unit, a new entry), and
  the number \(D\) of space dimensions (here \(D=3\) for the
  four-dimensional world and \(D=7\) for the eight-dimensional
  manifolds).
\item
  \textbf{New} (85 entries): every new symbol carries the suffix LL (for
  \emph{Lorentz super-loop}), written as a subscript (\(Q_{\rm LL}\)),
  or, when the lower position holds an index, as a superscript
  (\(\sigma^{\rm LL}_j\), \(L^{\rm LL}_{\mathbf x}\)). The exceptions
  are symbols that cannot be read as anything else: the bi-octonions and
  their parts in bold upright type (\(\mathbf x\), \(\mathbf e_p\),
  \(\mathbf u\), \(\mathbf b\), \ldots), the unit matrix \(\mathsf 1_4\)
  and the operators \(\mathsf O_1\), \(\mathsf O_2\), \(\mathsf O_3\),
  the number sets and algebras in blackboard type (\(\mathbb R\),
  \(\mathbb C\), \(\mathbb H\), \(\mathbb O\)), the groups and Lie
  algebras (\(SL(2,\mathbb C)\), \(\mathfrak{so}(3,1)\),
  \(\mathfrak g_2\), \ldots), operation signs, the labels of checks,
  theorems and hypotheses, and three plain symbols that no master
  registry uses with another meaning here: \(i\) (the imaginary unit),
  \(p\) (an index over the basis of an algebra) and \(\iota\) (the
  isospin index).
\end{itemize}

A superscript in upright type is a part of the name of a symbol, not a
symbol of its own: the ``F'' of \(k^{\rm F}_{\rm LL}\) and
\(\nu^{\rm F}_{\rm LL}\) means ``Fermi'', the ``F'' of
\(U(1)^{\rm F}_{\rm LL}\) and \(n^{\rm int,F}_{\rm LL}\) means
``fermion'', the ``E'' of \(k^{\rm E}_{\rm LL}\) means ``Euclidean'';
each whole symbol is registered once.

Each entry also lists the \textbf{plain-text names} under which the
programs, the figures and the data files write the symbol (for example
\texttt{Delta}, \texttt{mu}, \texttt{Lambda}, \texttt{dvec}). Every name
belongs to one entry (check LL-S1), and every column of every data file
is such a name or a descriptive word (check LL-S3).

\hypertarget{the-collisions-found-and-how-each-was-resolved}{%
\section{3. The collisions found, and how each was
resolved}\label{the-collisions-found-and-how-each-was-resolved}}

\begin{longtable}[]{@{}
  >{\raggedright\arraybackslash}p{(\columnwidth - 4\tabcolsep) * \real{0.3333}}
  >{\raggedright\arraybackslash}p{(\columnwidth - 4\tabcolsep) * \real{0.3333}}
  >{\raggedright\arraybackslash}p{(\columnwidth - 4\tabcolsep) * \real{0.3333}}@{}}
\toprule\noalign{}
\begin{minipage}[b]{\linewidth}\raggedright
collision
\end{minipage} & \begin{minipage}[b]{\linewidth}\raggedright
where
\end{minipage} & \begin{minipage}[b]{\linewidth}\raggedright
resolution
\end{minipage} \\
\midrule\noalign{}
\endhead
\bottomrule\noalign{}
\endlastfoot
the letter H: the quaternions, the Hamiltonian of the superfluid and the
Bogoliubov-de Gennes matrix (and, in the master registries, the Hubble
rates) & programs, data, text & quaternions \(\mathbb H\) (\texttt{H}),
Hamiltonian \(\mathcal H^{\rm sf}_{\rm LL}\) (\texttt{Hsf}),
Bogoliubov-de Gennes matrix \(\mathsf H_{\rm LL}\) (\texttt{H\_BdG}) \\
the letter phi: the automorphisms of the bi-octonions, the pair field of
the Yukawa completion, and the expectation value of the scalar field of
the main project & programs, text & automorphisms
\(\varphi^{\rm LL}_{\mathbf g,\mathbf v}\) (\texttt{varphi}), pair field
\(\phi_{\rm LL}\) (\texttt{phi}) \\
the letter d: the order parameter, the differential, and the gaps of the
main project & programs, data, text & order parameter
\(\vec d_{\rm LL}\) (\texttt{dvec}) with components \(d^\iota_{\rm LL}\)
(\texttt{d\_iota}); the letter d alone is the differential \\
the letter J: the bubble of the Yukawa completion (first draft), the
rotation generators, and the index sets and currents of the master
registries & program, text & rotations \(J^j_{\rm LL}\); bubble
\(I^{\rm bub}_{\rm LL}\) (\texttt{Ibub}) \\
the letter k: an index of the first drafts, the momenta, and the
wavenumbers of the master registries & program, text & isospin index
\(\iota\), space index \(j\); momenta \(\vec q_{\rm LL}\),
\(k^{\rm E}_{\rm LL}\), \(k^{\rm F}_{\rm LL}\) \\
the letter K: the boosts, the pair matrices, and the index sets and
evolution generator of the master registries & text & boosts
\(K^j_{\rm LL}\), pair matrices \(\mathsf K^\iota_{\rm LL}\)
(sans-serif, \texttt{Kch}) \\
the letter E: the energies, the matrix \(i\sigma^{\rm LL}_2\), and the
energies of the master registries & programs, text & energies
\(E^\pm_{\rm LL}\), \(\tilde E^\pm_{\rm LL}\); the matrix
\(\mathsf E_{\rm LL}\) (sans-serif, \texttt{eps2}) \\
the letter D: the number of space dimensions (master) and the
derivations & text & \(D\) alone is the number of space dimensions;
derivations \(\mathcal D_{\rm LL}\) (a general one),
\(D_{\rm LL}(\mathbf u)\), \(D^{\rm v}_{\rm LL}(\mathbf t)\) \\
the coordinates \texttt{x0}, \ldots, \texttt{x7} of the project's
spacetime and the coordinates of \(\mathbb R^{1,7}_{\rm LL}\) & program,
data, text & \(x^p_{\rm LL}\) (\texttt{x\_LL\^{}p}) \\
the Greek letters of the superfluid (in plain text \texttt{mu},
\texttt{Lambda}, \texttt{Omega}, \texttt{nu}, \texttt{Delta},
\texttt{zeta}, \texttt{eps}), which the master registries use for the
quantities of their own models & text & always with the suffix:
\(\mu_{\rm LL}\), \(\Lambda_{\rm LL}\), \(\Omega_{\rm LL}\),
\(\nu^{\rm F}_{\rm LL}\), \(\Delta_{\rm LL}\), \(\zeta_{\rm LL}\),
\(\varepsilon_{\rm LL}\) (check LL-S5) \\
the labels of the theorems of the superfluid, first \texttt{S1} to
\texttt{S4}, which are the labels of provenance documents of the main
project & text & SF1 to SF4 \\
letters of several objects that differ by type or suffix:
\(C_{\rm LL}\), \(\mathbb C\) and the checks C-1 to C-4;
\(L^{\rm LL}_{\mathbf x}\), \(\mathbb L_{\rm LL}\) and the checks L;
\(R^{\rm LL}_{\mathbf x}\), \(\mathbb R\), Theorem R and the checks R;
\(G_{\rm LL}\), \(\mathrm G_2\) and the checks G; \(N_{\rm LL}\) and the
checks N; \(Q_{\rm LL}\) and the checks Q; \(S^7_{\rm LL}\) and the
checks S; \(T^\iota_{\rm LL}\) and the transpose \({}^T\);
\(F_{\rm LL}\) and \(\mathcal F^{\rm A}_{\rm LL}\) & text & each is a
separate entry of the registry; a label of a check is always followed by
a hyphen and a number \\
plain-text names inside the meanings of the registry (\texttt{R\^{}8},
\texttt{M\_2(C)}, \ldots) that broke the LaTeX table & registry &
written as code; the LaTeX table sets them in typewriter type \\
the letter P: a generator in the registry of the main project, and the
first name of the sample point of the document on the \(3+1\) manifold &
program, data, figure & the words ``sample point'' \\
the letter G: the contact coupling \(G_{\rm LL}\) (\texttt{G}), and the
first name of the generator matrices in a report, with \texttt{R\_j} and
\texttt{B\_j} & program & the vector action
\(\mathcal V_{\rm LL}(\mathbf u)\) (\texttt{V\_LL}), an existing
entry \\
the name \texttt{chi}: the spinor \(\chi_{\rm LL}\), and an angle
variable of the quadrature of the bubble & program & the angle variable
is called \texttt{ang} \\
letters without suffix in the report, the comments and the figures of
the first version of the program \texttt{r13\_yukawa} (\texttt{X},
\texttt{M}, \texttt{x\^{}0}, \texttt{e\^{}}, the parameter
\(\theta_{\rm LL}\) written without its suffix, \texttt{k}) & program,
figures & the registered names \texttt{M\_LL(x)}, \texttt{x\_LL\^{}p},
\texttt{exp}, \texttt{theta}, \texttt{kE} \\
\end{longtable}

\hypertarget{symbols-shared-with-the-folder-revision_nogo}{%
\section{4. Symbols shared with the folder
Revision\_NoGo}\label{symbols-shared-with-the-folder-revision_nogo}}

The folder \texttt{Revision\_NoGo} of the repository has a registry of
its own, \texttt{provenance/symbols\_ng.csv}. The following symbols
occur in both registries and denote the same object in both folders:
\(i\) (the imaginary unit), \(\mathbb R\), \(\mathbb C\) (the real and
complex numbers), \(p\) (an index over the basis of an algebra),
\(\mathfrak g_2\), \(\mathrm G_2\) (the compact exceptional Lie algebra
and group), \(\mathfrak{so}(8)\), \(\mathfrak{so}(4,4)\),
\(\mathfrak{so}(8,\mathbb C)\), \(\mathfrak{so}(3,1)\),
\(\mathfrak{so}(3)\), \(\mathfrak{so}(4)\), \(\mathfrak{su}(2)\),
\(SU(2)\), \(U(1)\) (Lie algebras and groups), \(\mathrm{tr}\),
\(\det\), \(\mathrm{Re}\), \(\mathrm{diag}\), \({}^T\), \(\otimes\) and
\(\langle\ \rangle\) (operations). The script declares them in its list
\texttt{SHARED}, and check LL-S2 verifies that the shared symbols are
exactly the declared ones.

\hypertarget{plain-text-names-whose-meaning-needs-a-word}{%
\section{5. Plain-text names whose meaning needs a
word}\label{plain-text-names-whose-meaning-needs-a-word}}

The programs, the figures and the data files write formulas in plain
text, with the names listed in the last column of the table of Section
7. The following conventions complete that column.

\begin{itemize}
\tightlist
\item
  \textbf{Elements of the algebras}: \texttt{x}, \texttt{y}, \texttt{z},
  \texttt{w} are the bi-octonions \(\mathbf x\), \(\mathbf y\),
  \(\mathbf z\), \(\mathbf w\); \texttt{aH}, \texttt{bS}, \texttt{cH},
  \texttt{dS}, \texttt{bS\textquotesingle{}} are the bi-quaternions
  \(\mathbf a\), \(\mathbf b\), \(\mathbf c\), \(\mathbf d\),
  \(\mathbf b'\) of the form \(\mathbf a+\mathbf b\,\mathbf e_4\) (the
  letter H recalls that the element lies in
  \(\mathbb C\otimes\mathbb H\), the letter S that it multiplies
  \(\mathbf e_4\) to give a spinor charge); \texttt{uL},
  \texttt{uL\textquotesingle{}}, \texttt{uL1}, \texttt{uL2} are
  \(\mathbf u\), \(\mathbf u'\), \(\mathbf u_1\), \(\mathbf u_2\)
  (Lorentz part); \texttt{tI}, \texttt{tI\textquotesingle{}} are
  \(\mathbf t\), \(\mathbf t'\) (isospin); \texttt{gL}, \texttt{hL} are
  \(\mathbf g\), \(\mathbf h\) (elements of \(SL(2,\mathbb C)\));
  \texttt{vI} is \(\mathbf v\) (a unit quaternion, an isospin rotation);
  \texttt{rP}, \texttt{rP\textquotesingle{}} are \(\mathbf r\),
  \(\mathbf r'\) (translations). \texttt{e0}, \ldots, \texttt{e7} are
  the units \(\mathbf e_0,\dots,\mathbf e_7\) (and \texttt{1} in
  \texttt{gL\ =\ 1\ +\ i\ e1\ +\ e2} is \(\mathbf e_0\)); \texttt{ie1},
  \ldots, \texttt{ie7} are \(i\mathbf e_1,\dots,i\mathbf e_7\);
  \texttt{e\_p}, \texttt{e\_j}, \texttt{e\_iota} are \(\mathbf e_p\),
  \(\mathbf e_j\), \(\mathbf e_\iota\).
\item
  \textbf{Operations}: a final \texttt{bar} (\texttt{xbar},
  \texttt{bSbar}, \texttt{gLbar}) is the conjugation
  \(\bar{\mathbf x}\); \texttt{\^{}*} the complex conjugation;
  \texttt{\^{}dag} the Hermitian conjugation; \texttt{\^{}T} the
  transpose; \texttt{\^{}-1} the inverse; \texttt{(x)} the tensor
  product \(\otimes\); \texttt{o} the composition \(\circ\);
  \texttt{{[}x,y{]}} the commutator and \texttt{{[}x,y,z{]}} the
  associator; \texttt{Jac} the Jacobiator; \texttt{L\_x}, \texttt{R\_x}
  the left and right multiplications \(L^{\rm LL}_{\mathbf x}\),
  \(R^{\rm LL}_{\mathbf x}\), and \texttt{L\_{[}e1,e4{]}} the left
  multiplication by the commutator \([\mathbf e_1,\mathbf e_4]\);
  \texttt{cross(a,\ b)} is the cross product of two complex
  three-component vectors; \texttt{\textbar{}...\textbar{}} an absolute
  value or the length of a vector; the dot in \texttt{kE.(kE+pE)} the
  Euclidean scalar product; \texttt{x\ -\textgreater{}\ y} a map;
  \texttt{x(y)} a product or a function value, as the context of the
  formula shows.
\item
  \textbf{Coefficients of the families}: \texttt{c\_tw}, \texttt{c\_sl},
  \texttt{c\_PP}, \texttt{c\_Ps}, \texttt{c\_sP}, \texttt{c\_TT},
  \texttt{c\_Ts}, \texttt{c\_sT} are \(c^{\rm tw}_{\rm LL}\), \ldots,
  \(c^{\rm sT}_{\rm LL}\); a complex one is split as
  \texttt{ctw\_re\ +\ i\ ctw\_im} into its real and imaginary parts (and
  so for the others); \texttt{csP} and \texttt{cTT} are real. In checks
  M-3, M-4, M-7 and M-9 a failing triple or 4-tuple is given by the
  positions of its elements in the basis of the family: positions 0, 1,
  2 are \(\mathbf e_1\), \(\mathbf e_2\), \(\mathbf e_3\) and 3, 4, 5
  are \(i\mathbf e_1\), \(i\mathbf e_2\), \(i\mathbf e_3\) (the Lorentz
  part); in \(\mathcal F^{\rm B}_{\rm LL}\) positions 6 to 9 are the
  translations \(\mathbf e_0\), \(i\mathbf e_1\), \(i\mathbf e_2\),
  \(i\mathbf e_3\), and in \(\mathcal F^{\rm C}_{\rm LL}\) positions 6
  to 8 are the isospin elements \(\mathbf e_1\), \(\mathbf e_2\),
  \(\mathbf e_3\); the last eight positions are the spinor charges
  \(\mathbf e_4,\dots,\mathbf e_7\),
  \(i\mathbf e_4,\dots,i\mathbf e_7\).
\item
  \textbf{The superfluid}: \texttt{dvec} is \(\vec d_{\rm LL}\) and
  \texttt{d\_iota} its components; \texttt{eps2} is
  \(\mathsf E_{\rm LL}\), \texttt{eps3} the Levi-Civita symbol
  \(\epsilon^{\rm LL}_{jj'j''}\); \texttt{E\^{}(+-)} stands for
  \(E^+_{\rm LL}\) and \(E^-_{\rm LL}\), and \texttt{E\^{}(+-)(Delta)}
  for these energies with the gap \texttt{Delta} in place of
  \(\Delta_{\rm LL}\); \texttt{Delta\_(+-)} are the gaps
  \(\Delta^\pm_{\rm LL}\); \texttt{Etilde\^{}(+-)} are
  \(\tilde E^\pm_{\rm LL}\); \texttt{eps} is \(\varepsilon_{\rm LL}\),
  \texttt{q} the momentum \(\vec q_{\rm LL}\) or its length;
  \texttt{Q\_LL(Kch)} is the electric charge of the two spinor charges
  that the pair operator with the matrix \texttt{Kch} annihilates;
  \texttt{0+} and \texttt{0-} are spin zero with positive and negative
  parity; \texttt{spin(spin+1)} in check Q-3 is the spin quantum number
  times itself plus one, the value 0, 2 or \(\tfrac34\) of the Casimir
  for the spins 0, 1 and \(\tfrac12\). In the folder
  \texttt{explanations}: \texttt{Dpair} is the isospin part
  \(\mathsf D_{\rm LL}\) of the pairing block (not \texttt{D\_LL}, which
  is the Lorentz derivation); \texttt{F\_LL/2\ -\ T\_LL\^{}3} and
  \texttt{F\_LL/2\ +\ T\_LL\^{}3} are
  \(\tfrac12F_{\rm LL}-T^3_{\rm LL}\) and
  \(\tfrac12F_{\rm LL}+T^3_{\rm LL}\), the numbers of the isospin-down
  and of the isospin-up particles; \texttt{fraction(+1)} is the fraction
  of the condensed pairs of the charge \(+1\). In the program
  \texttt{r13\_yukawa} of the folder \texttt{explanations}:
  \texttt{Lmat} is \(\mathsf L_{\rm LL}(\mathbf g)\) and
  \texttt{Lmat{[}0{]}{[}0{]}} its time-time entry; \texttt{V\_LL(uL)} is
  the vector action \(\mathcal V_{\rm LL}(\mathbf u)\), for example
  \texttt{V\_LL(J\_LL\^{}j)}; \texttt{chi} is \(\chi_{\rm LL}\);
  \texttt{theta} is \(\theta_{\rm LL}\) and \texttt{exp(theta)} its
  exponential; \texttt{xi\_p} is \(\xi^{\rm LL}_p\);
  \texttt{eps3\_jj\textquotesingle{}j\textquotesingle{}\textquotesingle{}}
  is \(\epsilon^{\rm LL}_{jj'j''}\) and
  \texttt{sum\_j\textquotesingle{}\textquotesingle{}} a sum over
  \(j''\); \texttt{{[}{[}a,\ b{]},\ {[}c,\ d{]}{]}} with numbers in
  place of the letters is a \(2\times2\) matrix written row by row;
  \texttt{sqrt3} is \(\sqrt3\);
  \texttt{int\_(\textbar{}kE\textbar{}\ \textless{}\ Lambda)\ d\^{}4kE}
  is the integral over the four-dimensional ball of radius
  \(\Lambda_{\rm LL}\); \texttt{dIbub/d(pE\^{}2)} is the derivative of
  \(I^{\rm bub}_{\rm LL}\) with respect to \((p^{\rm E}_{\rm LL})^2\);
  \texttt{1e-12} is \(10^{-12}\); ``the sample point'' is the point with
  the coordinates \((2,1.2,1.5,1.0)\).
\item
  \textbf{Abbreviations}: \texttt{BdG} is Bogoliubov-de Gennes,
  \texttt{BCS} Bardeen-Cooper-Schrieffer; \texttt{M\_2(C)} the complex
  \(2\times2\) matrices; \texttt{R\^{}8}, \texttt{C\^{}3},
  \texttt{C\^{}8} the real or complex vector spaces of that dimension;
  \texttt{Im\ H}, \texttt{Im\ O} and \texttt{C\ (x)\ H} are
  \({\rm Im}\,\mathbb H\), \({\rm Im}\,\mathbb O\) and
  \(\mathbb C\otimes\mathbb H\); the groups and algebras
  (\texttt{SL(2,C)}, \texttt{so(3,1)}, \texttt{g2(C)},
  \texttt{SU(2)\_v}, \texttt{U(1)\_F\_LL}, \texttt{Z2}, \ldots) are
  written as in the text; \texttt{4x4} means four by four; ``exact'' in
  a check means integer or rational arithmetic, ``double precision''
  floating-point arithmetic with 53-bit mantissas.
\item
  \textbf{Comments of the source code} also call the quaternion units
  \(\mathbf e_1\), \(\mathbf e_2\), \(\mathbf e_3\) \texttt{i\_q},
  \texttt{j\_q}, \texttt{k\_q}, and the unit \(\mathbf e_4\) of the
  Cayley-Dickson construction \texttt{l}; these names do not occur in
  the reports, the data files or the figures.
\end{itemize}

\hypertarget{the-machine-checks-of-the-symbols}{%
\section{6. The machine checks of the
symbols}\label{the-machine-checks-of-the-symbols}}

The script \texttt{code/scripts/check\_symbols.py} verifies eight
statements:

\begin{itemize}
\tightlist
\item
  \textbf{LL-S1} no symbol text is registered twice, and no plain-text
  name belongs to two entries;
\item
  \textbf{LL-S2} every taken entry has exactly the meaning of the master
  registry from which it is taken; no item of a new entry is an item of
  either master registry (an item is one formula of a symbol text, for
  example \(\mathbf x\) in ``\(\mathbf x\), \(\mathbf y\)''); and the
  items that the registry of \texttt{Revision\_NoGo} also contains are
  exactly the declared shared ones (Section 4);
\item
  \textbf{LL-S3} every column of every data file \texttt{data/sl\_*.csv}
  and \texttt{explanations/data/ex\_*.csv} is a registered plain-text
  name or a descriptive word, and every data file states its content in
  a header line \texttt{\#\ content:};
\item
  \textbf{LL-S4} every document of the folder (and the manuscript)
  contains the complete table of symbols;
\item
  \textbf{LL-S5} in the formulas of the documents, the symbols of the
  superfluid \(\mu_{\rm LL}\), \(\Lambda_{\rm LL}\),
  \(\Omega_{\rm LL}\), \(\nu^{\rm F}_{\rm LL}\), \(\Delta_{\rm LL}\),
  \(\zeta_{\rm LL}\), \(\varepsilon_{\rm LL}\) always carry the suffix
  LL (a superscript may stand before it);
\item
  \textbf{LL-S6} in the formulas of the documents and in the meanings of
  the new entries, every Latin letter other than the index \(j\), \(D\),
  the imaginary unit \(i\), the index \(p\) and the differential d
  carries a script that contains the suffix LL, and the Greek letters
  that are symbols carry it as well (except the index \(\iota\) and the
  number \(\pi\)); letters in bold, sans-serif, calligraphic, blackboard
  and fraktur type are checked by LL-S7;
\item
  \textbf{LL-S7} every letter in bold, sans-serif, calligraphic,
  blackboard or fraktur type that occurs in a formula of a document is
  registered;
\item
  \textbf{LL-S8} every name with the suffix \texttt{\_LL}, every name of
  an algebra element (\texttt{uL}, \texttt{tI}, \texttt{gL},
  \texttt{hL}, \texttt{vI}, \texttt{aH}, \texttt{bS}, \texttt{cH},
  \texttt{dS}, \texttt{rP} and their variants), every coefficient
  (\texttt{c\_tw}, \ldots) and every pair matrix (\texttt{Kch}\ldots)
  that occurs in a report or a data file of the folder
  (\texttt{data/sl\_*} and \texttt{explanations/data/ex\_*}) is a
  registered plain-text name.
\end{itemize}

A line break inside a formula counts as a space, as for LaTeX. The
output of the last run:

\begin{verbatim}
PROJECT spacetimes, Revision_SuperLoop: checks of the symbols (registry provenance/symbols_ll.csv)
8 of 8 checks passed

[PASS] LL-S1  every symbol text is registered once (87 entries) and every plain-text name of the programs, figures and data files belongs to one entry (310 names)
[PASS] LL-S2  the 2 taken entries carry exactly the meaning of the MASTER registry from which they are taken; none of the 203 items of the 85 new entries is an item of either master registry (503 items compared); the 22 items that the registry of Revision_NoGo also contains are exactly the 22 declared shared ones, each with the same object ($SU(2)$, $U(1)$, $\det$, $\langle\ \rangle$, $\mathbb C$, $\mathbb R$, $\mathfrak g_2$, $\mathfrak{so}(3)$, $\mathfrak{so}(3,1)$, $\mathfrak{so}(4)$, $\mathfrak{so}(4,4)$, $\mathfrak{so}(8)$, $\mathfrak{so}(8,\mathbb C)$, $\mathfrak{su}(2)$, $\mathrm G_2$, $\mathrm{Re}$, $\mathrm{diag}$, $\mathrm{tr}$, $\otimes$, $i$, $p$, ${}^T$)
[PASS] LL-S3  every column of the 23 data files data/sl_*.csv and explanations/data/ex_*.csv is a registered plain-text name or a descriptive word (120 different columns), and every file states its content in a line '# content:'
[PASS] LL-S4  each of the 6 documents of Revision_SuperLoop and the manuscript contains the complete table of the 87 registered symbols, so that every symbol it uses is defined in the document itself
[PASS] LL-S5  in the 3098 formulas of the documents, mu, Lambda, Omega, nu, Delta, zeta and varepsilon always carry the suffix LL
[PASS] LL-S6  in the 3620 formulas of the documents and of the new registry entries, no Latin letter except the index j, D, the imaginary unit i, the index p and the differential d stands without the suffix LL in its scripts, and the Greek letters that are symbols carry the suffix as well
[PASS] LL-S7  every one of the 3010 occurrences of a letter in bold, sans-serif, calligraphic, blackboard or fraktur type in the formulas of the documents is a registered symbol (47 registered styled letters)
[PASS] LL-S8  every one of the 3158 occurrences of a name with the suffix _LL, of an algebra element, of a family coefficient or of a pair matrix in the 32 reports and data files (data/sl_* and explanations/data/ex_*) is a registered plain-text name
\end{verbatim}

\hypertarget{the-complete-table-of-the-symbols}{%
\section{7. The complete table of the
symbols}\label{the-complete-table-of-the-symbols}}

\textbf{1 Indices, conventions and number sets}

\begin{longtable}[]{@{}
  >{\raggedright\arraybackslash}p{(\columnwidth - 6\tabcolsep) * \real{0.2200}}
  >{\raggedright\arraybackslash}p{(\columnwidth - 6\tabcolsep) * \real{0.5000}}
  >{\raggedright\arraybackslash}p{(\columnwidth - 6\tabcolsep) * \real{0.0800}}
  >{\raggedright\arraybackslash}p{(\columnwidth - 6\tabcolsep) * \real{0.2000}}@{}}
\toprule\noalign{}
\begin{minipage}[b]{\linewidth}\raggedright
symbol
\end{minipage} & \begin{minipage}[b]{\linewidth}\raggedright
meaning
\end{minipage} & \begin{minipage}[b]{\linewidth}\raggedright
origin
\end{minipage} & \begin{minipage}[b]{\linewidth}\raggedright
plain-text names
\end{minipage} \\
\midrule\noalign{}
\endhead
\bottomrule\noalign{}
\endlastfoot
\(i,j\) & index of space-like coordinates; values 1, 2, 3 (ordinary
space) unless a sum is said to run over all space-like coordinates of
the spacetime. In Revision\_SuperLoop only \(j\) is used as an index,
with the values 1, 2, 3, and with primes (\(j'\), \(j''\)) where two or
three such indices are needed: it labels the three imaginary quaternion
units \(\mathbf e_j\), which correspond to the three space directions of
\(\mathbb R^{1,3}_{\rm LL}\), the coordinates \(x^j_{\rm LL}\), the
momentum components \(q^j_{\rm LL}\), the rotation and boost generators
\(J^j_{\rm LL}\), \(K^j_{\rm LL}\), and the Pauli matrices
\(\sigma^{\rm LL}_j\) and Dirac matrices \(\alpha^j_{\rm LL}\),
\(\gamma^j_{\rm LL}\). The letter \(i\) is the imaginary unit (entry
\(i\)); the isospin index is \(\iota\) (a different letter). & taken &
\texttt{j} \texttt{j\textquotesingle{}}
\texttt{j\textquotesingle{}\textquotesingle{}} \\
\(D\) & the number of space dimensions in which the particles of matter
have momentum: 7 (all space coordinates of 7+1), 4 (fields that do not
depend on \(x5,x6,x7\)) or 3 (fields that depend on ordinary space
only). \(V_D=\pi^{D/2}/\Gamma_{\rm E}(D/2+1)\) is the volume and
\(A_D=D\,V_D\) the area of the unit ball in \(D\) dimensions
(\(\Gamma_{\rm E}\): Euler's gamma function). In Revision\_SuperLoop
\(D=3\) for the four-dimensional world of the superfluid and \(D=7\) for
the eight-dimensional manifolds (on \(4+4\) power counting sees the
total dimension 8 as well); the value 4, the coordinates \(x5,x6,x7\) of
the eight-dimensional spacetime of the project (on which such fields do
not depend), and the volume and the area of the unit ball do not occur.
& taken & \texttt{D} \\
\(i\) & the imaginary unit, \(i^2=-1\); in \(\mathbb B_{\rm LL}\) it is
the imaginary unit of the complex coefficients: it commutes and
associates with every bi-octonion and is not one of the units
\(\mathbf e_1,\dots,\mathbf e_7\) & new & \\
\(\mathbb R\), \(\mathbb C\) & the real and the complex numbers (plain
text: \texttt{R\^{}8}, \texttt{C\^{}3}, \texttt{C\^{}8} are the real or
complex vector spaces of that dimension) & new & \texttt{R} \texttt{C}
\texttt{R\^{}8} \texttt{C\^{}3} \texttt{C\^{}8} \\
\(p\), \(p'\) & a summation index over the basis of an algebra (values
\(0,\dots,7\), or \(0,\dots,3\) for the quaternions): in
Revision\_SuperLoop it labels the units \(\mathbf e_p\) of
\(\mathbb B_{\rm LL}\), the coefficients \(\xi^{\rm LL}_p(\mathbf x)\),
the coordinates \(x^p_{\rm LL}\) of \(\mathbb R^{1,7}_{\rm LL}\) and
\(\mathbb R^{1,3}_{\rm LL}\) and the Dirac matrices
\(\gamma^p_{\rm LL}\) (\(p=0,1,2,3\)); \(p'\) is a second index of the
same kind (plain text \texttt{p1}, \texttt{p2}, \texttt{p3}: three such
indices) & new & \texttt{p} \texttt{p\textquotesingle{}} \texttt{p1}
\texttt{p2} \texttt{p3} \\
\(\iota\), \(\iota'\), \(\iota''\) & the isospin index, values 1, 2, 3:
it labels the isospin generators \(T^\iota_{\rm LL}\), the Pauli
matrices \(\tau^{\rm LL}_\iota\) on the isospin index of a field, the
pair channels \(\mathsf K^\iota_{\rm LL}\),
\(\mathsf K'^\iota_{\rm LL}\), the pair operators \(\Pi^\iota_{\rm LL}\)
and the components \(d^\iota_{\rm LL}\) of the order parameter; it is an
internal direction, not a direction in space (that is \(j\)) & new &
\texttt{iota} \\
\(\epsilon^{\rm LL}_{jj'j''}\) & the Levi-Civita symbol: \(+1\) if
\((j,j',j'')\) is an even permutation of \((1,2,3)\), \(-1\) if it is an
odd one, \(0\) otherwise; also with the isospin indices,
\(\epsilon^{\rm LL}_{\iota\iota'\iota''}\) & new & \texttt{eps3} \\
\(\sigma^{\rm LL}_0\), \(\sigma^{\rm LL}_j\) & the \(2\times2\) unit
matrix \(\sigma^{\rm LL}_0\) and the Pauli matrices
\(\sigma^{\rm LL}_1=\begin{pmatrix}0&1\\1&0\end{pmatrix}\),
\(\sigma^{\rm LL}_2=\begin{pmatrix}0&-i\\i&0\end{pmatrix}\),
\(\sigma^{\rm LL}_3=\begin{pmatrix}1&0\\0&-1\end{pmatrix}\), acting on
the two components of a Weyl spinor or on the species index & new &
\texttt{sigma\_0} \texttt{sigma\_j} \texttt{sigma\_1} \texttt{sigma\_2}
\texttt{sigma\_3} \\
\(\tau^{\rm LL}_\iota\) & the same three Pauli matrices acting on the
isospin index of a field (the index of the two components up and down of
an isodoublet) & new & \texttt{tau\_iota} \texttt{tau\_1}
\texttt{tau\_2} \texttt{tau\_3} \\
\(\mathsf E_{\rm LL}\) & the real antisymmetric \(2\times2\) matrix
\(i\sigma^{\rm LL}_2=\begin{pmatrix}0&1\\-1&0\end{pmatrix}\), with
\(\mathsf E_{\rm LL}^{-1}=\mathsf E_{\rm LL}^T=-\mathsf E_{\rm LL}\); it
acts on a Weyl, isospin or species index, as stated at each place & new
& \texttt{eps2} \\
\(\mathsf 1_4\) & the \(4\times4\) unit matrix & new & \\
\end{longtable}

\textbf{2 The bi-octonions and the super-loop}

\begin{longtable}[]{@{}
  >{\raggedright\arraybackslash}p{(\columnwidth - 6\tabcolsep) * \real{0.2200}}
  >{\raggedright\arraybackslash}p{(\columnwidth - 6\tabcolsep) * \real{0.5000}}
  >{\raggedright\arraybackslash}p{(\columnwidth - 6\tabcolsep) * \real{0.0800}}
  >{\raggedright\arraybackslash}p{(\columnwidth - 6\tabcolsep) * \real{0.2000}}@{}}
\toprule\noalign{}
\begin{minipage}[b]{\linewidth}\raggedright
symbol
\end{minipage} & \begin{minipage}[b]{\linewidth}\raggedright
meaning
\end{minipage} & \begin{minipage}[b]{\linewidth}\raggedright
origin
\end{minipage} & \begin{minipage}[b]{\linewidth}\raggedright
plain-text names
\end{minipage} \\
\midrule\noalign{}
\endhead
\bottomrule\noalign{}
\endlastfoot
\(\mathbb H\), \(\mathbb O\) & the quaternions and the octonions: the
real combinations of \(\mathbf e_0,\dots,\mathbf e_3\) and of
\(\mathbf e_0,\dots,\mathbf e_7\); \({\rm Im}\,\mathbb H\) and
\({\rm Im}\,\mathbb O\) are their imaginary parts, the real combinations
of \(\mathbf e_1,\mathbf e_2,\mathbf e_3\) and of
\(\mathbf e_1,\dots,\mathbf e_7\) (plain text \texttt{Im\ H},
\texttt{Im\ O}; \texttt{C\ (x)\ H} is \(\mathbb C\otimes\mathbb H\), the
bi-quaternions) & new & \texttt{H} \texttt{O} \\
\(\mathbb B_{\rm LL}\) & the bi-octonions \(\mathbb C\otimes\mathbb O\):
the complex combinations
\(\mathbf x=\sum_{p=0}^7\xi^{\rm LL}_p(\mathbf x)\,\mathbf e_p\) of the
eight units, with the product of the units extended bilinearly; real
dimension 16 & new & \texttt{B\_LL} \\
\(\mathbf e_0,\dots,\mathbf e_7\); \(\mathbf e_p\), \(\mathbf e_j\),
\(\mathbf e_\iota\) & the units of \(\mathbb B_{\rm LL}\), of
\(\mathbb O\) and of \(\mathbb H\): \(\mathbf e_0=1\);
\(\mathbf e_1,\mathbf e_2,\mathbf e_3\) the quaternion units,
\(\mathbf e_j^2=-1\),
\(\mathbf e_1\mathbf e_2=-\mathbf e_2\mathbf e_1=\mathbf e_3\) and
cyclically; \(\mathbf e_4\) the unit of the Cayley-Dickson doubling,
\(\mathbf e_4^2=-1\); \(\mathbf e_5=\mathbf e_1\mathbf e_4\),
\(\mathbf e_6=\mathbf e_2\mathbf e_4\),
\(\mathbf e_7=\mathbf e_3\mathbf e_4\) (\(\mathbf e_\iota\): a
quaternion unit labelled by the isospin index; plain text \texttt{ie1},
\ldots, \texttt{ie7}: the bi-octonions
\(i\mathbf e_1,\dots,i\mathbf e_7\)) & new & \texttt{e0} \texttt{e1}
\texttt{e2} \texttt{e3} \texttt{e4} \texttt{e5} \texttt{e6} \texttt{e7}
\texttt{e\_p} \texttt{e\_p\textquotesingle{}} \texttt{e\_j}
\texttt{e\_iota} \texttt{e\_p1} \texttt{e\_p2} \texttt{e\_p3}
\texttt{ie1} \texttt{ie2} \texttt{ie3} \texttt{ie4} \texttt{ie5}
\texttt{ie6} \texttt{ie7} \\
\(\xi^{\rm LL}_p(\mathbf x)\) & the complex coefficient of the unit
\(\mathbf e_p\) in the bi-octonion \(\mathbf x\) & new & \texttt{xi\_p}
\texttt{xi\_0} \texttt{xi\_1} \texttt{xi\_2} \texttt{xi\_3} \\
\(\mathbf x\), \(\mathbf y\), \(\mathbf z\), \(\mathbf w\) & elements of
\(\mathbb B_{\rm LL}\) (bi-octonions); \(\mathbf w\) is also a point of
the super-loop manifold \(\mathbb L_{\rm LL}\) & new & \texttt{x}
\texttt{y} \texttt{z} \texttt{w} \\
\(\mathbf a\), \(\mathbf b\), \(\mathbf c\), \(\mathbf d\);
\(\mathbf b'\) & bi-quaternions: elements of
\(\mathbb C\otimes\mathbb H\), the complex combinations of
\(\mathbf e_0,\mathbf e_1,\mathbf e_2,\mathbf e_3\); every bi-octonion
is written \(\mathbf x=\mathbf a+\mathbf b\,\mathbf e_4\)
(Cayley-Dickson form), and \(\mathbf b\,\mathbf e_4\),
\(\mathbf b'\,\mathbf e_4\) are spinor charges & new & \texttt{aH}
\texttt{bS} \texttt{cH} \texttt{dS} \texttt{bS\textquotesingle{}} \\
\(\bar{\mathbf x}\), \(\mathbf x^*\), \(\mathbf x^\dagger\) & the three
conjugations (also of bi-quaternions): \(\bar{\mathbf x}\) changes the
signs of the coefficients of \(\mathbf e_1,\dots,\mathbf e_7\) (it is
complex-linear,
\(\overline{\mathbf x\mathbf y}=\bar{\mathbf y}\,\bar{\mathbf x}\), and
\(\overline{\mathbf a+\mathbf b\mathbf e_4}=\bar{\mathbf a}-\mathbf b\mathbf e_4\));
\(\mathbf x^*\) replaces every complex coefficient by its complex
conjugate; \(\mathbf x^\dagger=(\bar{\mathbf x})^*\) (plain text:
\texttt{xbar}, \texttt{x\^{}*}, \texttt{x\^{}dag}, and likewise
\texttt{aHbar}, \texttt{gL\^{}dag}, \ldots) & new & \texttt{xbar}
\texttt{x\^{}*} \texttt{x\^{}dag} \texttt{aHbar} \texttt{bSbar}
\texttt{bS\textquotesingle{}bar} \texttt{cHbar} \texttt{dSbar}
\texttt{gLbar} \texttt{hLbar} \texttt{rPbar}
\texttt{rP\textquotesingle{}bar} \texttt{bS\^{}*}
\texttt{bS\textquotesingle{}\^{}*} \texttt{gL\^{}dag}
\texttt{uL\^{}dag} \\
\(N_{\rm LL}(\mathbf x)\), \(N_{\rm LL}(\mathbf x,\mathbf y)\) & the
complex norm
\(N_{\rm LL}(\mathbf x)=\sum_p\xi^{\rm LL}_p(\mathbf x)^2\), with
\(\mathbf x\bar{\mathbf x}=\bar{\mathbf x}\mathbf x=N_{\rm LL}(\mathbf x)\,\mathbf e_0\),
and its polarization
\(N_{\rm LL}(\mathbf x,\mathbf y)=\sum_p\xi^{\rm LL}_p(\mathbf x)\,\xi^{\rm LL}_p(\mathbf y)\),
\(N_{\rm LL}(\mathbf x,\mathbf x)=N_{\rm LL}(\mathbf x)\)
(complex-valued; for bi-quaternions the determinant of the \(2\times2\)
matrix \(M_{\rm LL}\)) & new & \texttt{N\_LL} \\
\([\mathbf x,\mathbf y]\); \([\mathbf x,\mathbf y,\mathbf z]\) & the
commutator \(\mathbf x\mathbf y-\mathbf y\mathbf x\) and the associator
\((\mathbf x\mathbf y)\mathbf z-\mathbf x(\mathbf y\mathbf z)\) of
bi-octonions (the same brackets of other elements are defined at their
place: the bracket of an algebra \(\mathcal F^{\rm A}_{\rm LL}\),
\(\mathcal F^{\rm B}_{\rm LL}\), \(\mathcal F^{\rm C}_{\rm LL}\), and
the commutator of two matrices or operators, the product in one order
minus the product in the other) & new & \texttt{{[}x,y{]}}
\texttt{{[}x,y,z{]}} \\
\(\mathrm{Jac}_{\rm LL}\) & the Jacobiator
\(\mathrm{Jac}_{\rm LL}(\mathbf x,\mathbf y,\mathbf z)=[[\mathbf x,\mathbf y],\mathbf z]+[[\mathbf y,\mathbf z],\mathbf x]+[[\mathbf z,\mathbf x],\mathbf y]\)
of a bracket; it vanishes identically exactly in a Lie algebra & new &
\texttt{Jac} \\
\(L^{\rm LL}_{\mathbf x}\), \(R^{\rm LL}_{\mathbf x}\) & the left and
right multiplications \(\mathbf y\mapsto\mathbf x\mathbf y\) and
\(\mathbf y\mapsto\mathbf y\mathbf x\) by a bi-octonion, real
\(16\times16\) matrices (on \(\mathbb O\) or
\(\mathbb O^{\rm s}_{\rm LL}\): \(8\times8\)) & new & \texttt{L\_x}
\texttt{L\_y} \texttt{L\_e1} \texttt{L\_e4} \texttt{L\_{[}e1,e4{]}}
\texttt{L\_(e\_p)} \texttt{L\_(e\_p\textquotesingle{})} \texttt{R\_x}
\texttt{R\_y} \\
\(\mathbb L_{\rm LL}\) & the Lorentz super-loop: the unit bi-octonions,
\(\mathbb L_{\rm LL}=\{\mathbf w\in\mathbb B_{\rm LL}:N_{\rm LL}(\mathbf w)=1\}\),
with the product of \(\mathbb B_{\rm LL}\); a Moufang loop (not a
group), and a manifold of real dimension 14 & new & \texttt{L\_LL} \\
\(\mathbf g\), \(\mathbf h\) & elements of \(SL(2,\mathbb C)\),
i.e.~unit bi-quaternions (\(N_{\rm LL}(\mathbf g)=1\)), a subgroup of
\(\mathbb L_{\rm LL}\) & new & \texttt{gL} \texttt{hL} \\
\(\mathbf v\) & a unit (bi-)quaternion that rotates the isospin of the
spinor charges (from the left) & new & \texttt{vI} \\
\(\mathfrak m_{\rm LL}\), \(\mathfrak l_{\rm LL}\),
\(\mathfrak s_{\rm LL}\) & the Malcev algebra of the super-loop,
\(\mathfrak m_{\rm LL}=\mathbb C\otimes{\rm Im}\,\mathbb O\) with the
commutator (real dimension 14), its Lorentz part
\(\mathfrak l_{\rm LL}=\mathbb C\otimes{\rm Im}\,\mathbb H\) (6, a Lie
algebra isomorphic to \(\mathfrak{so}(3,1)\)), and the space of spinor
charges
\(\mathfrak s_{\rm LL}=(\mathbb C\otimes\mathbb H)\,\mathbf e_4\) (8);
\(\mathfrak m_{\rm LL}=\mathfrak l_{\rm LL}\oplus\mathfrak s_{\rm LL}\)
& new & \texttt{m\_LL} \texttt{l\_LL} \texttt{s\_LL} \\
\(J^j_{\rm LL}\), \(K^j_{\rm LL}\) & the rotation and boost generators
of \(\mathfrak l_{\rm LL}\), \(J^j_{\rm LL}=\mathbf e_j/2\) and
\(K^j_{\rm LL}=i\,\mathbf e_j/2\) & new & \texttt{J\_LL\^{}j}
\texttt{K\_LL\^{}j} \texttt{J\_LL\^{}j\textquotesingle{}}
\texttt{K\_LL\^{}j\textquotesingle{}}
\texttt{J\_LL\^{}j\textquotesingle{}\textquotesingle{}}
\texttt{K\_LL\^{}j\textquotesingle{}\textquotesingle{}}
\texttt{J\_LL\^{}3} \\
\(\mathbf u\), \(\mathbf u'\), \(\mathbf u_1\), \(\mathbf u_2\) &
elements of the Lorentz part \(\mathfrak l_{\rm LL}\) (Lorentz
generators): imaginary bi-quaternions & new & \texttt{uL}
\texttt{uL\textquotesingle{}} \texttt{uL1} \texttt{uL2} \\
\(\mathbf t\), \(\mathbf t'\) & elements of \({\rm Im}\,\mathbb H\)
(real imaginary quaternions): isospin generators, which act on the
spinor charges from the left & new & \texttt{tI}
\texttt{tI\textquotesingle{}} \\
\(W_{\rm LL}(\mathbf u)\) & the Weyl representation of
\(\mathfrak l_{\rm LL}\) on bi-quaternions,
\(W_{\rm LL}(\mathbf u)\mathbf b=-\mathbf b\,\mathbf u\): a Lie
representation,
\([W_{\rm LL}(\mathbf u_1),W_{\rm LL}(\mathbf u_2)]=W_{\rm LL}([\mathbf u_1,\mathbf u_2])\)
& new & \texttt{W\_LL} \\
\(\rho_{\rm LL}(\mathbf u)\) & the inner action of a Lorentz element on
the spinor charges,
\([\mathbf u,\mathbf b\,\mathbf e_4]=(\rho_{\rm LL}(\mathbf u)\mathbf b)\,\mathbf e_4\)
with \(\rho_{\rm LL}=-2W_{\rm LL}\); not a representation & new &
\texttt{rho\_LL} \\
\(\Theta_{\rm LL}\) & the twist charge: the linear map of
\(\mathfrak m_{\rm LL}\) that is \(1\) on \(\mathfrak l_{\rm LL}\) and
\(-2\) on \(\mathfrak s_{\rm LL}\); the inner action is
\([\mathbf u,\cdot\,]=D_{\rm LL}(\mathbf u)\circ\Theta_{\rm LL}\) (as a
number: the eigenvalue \(1\) or \(-2\), or the value of
\(c^{\rm tw}_{\rm LL}\) in the families) & new & \texttt{Theta\_LL}
\texttt{Theta\_LL\^{}2} \\
\(D_{\rm LL}(\mathbf u)\), \(D^{\rm v}_{\rm LL}(\mathbf t)\) & the
Lorentz and isospin derivations of \(\mathbb B_{\rm LL}\):
\(D_{\rm LL}(\mathbf u)(\mathbf a+\mathbf b\mathbf e_4)=[\mathbf u,\mathbf a]+(-\mathbf b\mathbf u)\mathbf e_4\)
and
\(D^{\rm v}_{\rm LL}(\mathbf t)(\mathbf a+\mathbf b\mathbf e_4)=(\mathbf t\mathbf b)\mathbf e_4\)
(both linear in their argument) & new & \texttt{D\_LL}
\texttt{Dv\_LL} \\
\(\mathcal D_{\rm LL}\) & a derivation of \(\mathbb B_{\rm LL}\): a
linear map with
\(\mathcal D_{\rm LL}(\mathbf x\mathbf y)=(\mathcal D_{\rm LL}\mathbf x)\mathbf y+\mathbf x(\mathcal D_{\rm LL}\mathbf y)\),
the infinitesimal form of an automorphism & new & \\
\(\varphi^{\rm LL}_{\mathbf g,\mathbf v}\) & the automorphism
\(\mathbf a+\mathbf b\mathbf e_4\mapsto\mathbf g\mathbf a\bar{\mathbf g}+(\mathbf v\mathbf b\bar{\mathbf g})\mathbf e_4\)
of \(\mathbb B_{\rm LL}\) (\(\mathbf g\) a Lorentz transformation,
\(\mathbf v\) an isospin rotation) & new & \texttt{varphi} \\
\(\varrho_{\rm LL}\); \(\mathsf O_1\), \(\mathsf O_2\), \(\mathsf O_3\)
& a linear map from \(\mathfrak m_{\rm LL}\) to the linear operators
(matrices) on a vector space that turns brackets into commutators,
\(\varrho_{\rm LL}([\mathbf x,\mathbf y])=\varrho_{\rm LL}(\mathbf x)\varrho_{\rm LL}(\mathbf y)-\varrho_{\rm LL}(\mathbf y)\varrho_{\rm LL}(\mathbf x)\)
(Theorem R); \(\mathsf O_1,\mathsf O_2,\mathsf O_3\) three such
operators & new & \\
\(\mathrm{Kil}_{\rm LL}\); \(\mathrm{ad}\) & the Killing form
\(\mathrm{Kil}_{\rm LL}(\mathbf x,\mathbf y)=\mathrm{tr}(\mathrm{ad}\,\mathbf x\,\mathrm{ad}\,\mathbf y)\)
of a Lie algebra, where \(\mathrm{ad}\,\mathbf x\) is the matrix of the
map \(\mathbf y\mapsto[\mathbf x,\mathbf y]\) & new & \texttt{Killing}
\texttt{ad} \\
\end{longtable}

\textbf{3 Operators on the super-loop and their eigenvalues}

\begin{longtable}[]{@{}
  >{\raggedright\arraybackslash}p{(\columnwidth - 6\tabcolsep) * \real{0.2200}}
  >{\raggedright\arraybackslash}p{(\columnwidth - 6\tabcolsep) * \real{0.5000}}
  >{\raggedright\arraybackslash}p{(\columnwidth - 6\tabcolsep) * \real{0.0800}}
  >{\raggedright\arraybackslash}p{(\columnwidth - 6\tabcolsep) * \real{0.2000}}@{}}
\toprule\noalign{}
\begin{minipage}[b]{\linewidth}\raggedright
symbol
\end{minipage} & \begin{minipage}[b]{\linewidth}\raggedright
meaning
\end{minipage} & \begin{minipage}[b]{\linewidth}\raggedright
origin
\end{minipage} & \begin{minipage}[b]{\linewidth}\raggedright
plain-text names
\end{minipage} \\
\midrule\noalign{}
\endhead
\bottomrule\noalign{}
\endlastfoot
\(Q_{\rm LL}\) & the electric charge, \(Q_{\rm LL}=T^3_{\rm LL}\):
eigenvalues \(0\) on \(\mathbb C\otimes\mathbb H\) and \(\pm\tfrac12\)
on the spinor charges & new & \texttt{Q\_LL} \\
\(T^\iota_{\rm LL}\), \((T_{\rm LL})^2\) & the isospin generators
\(T^\iota_{\rm LL}=i\,D^{\rm v}_{\rm LL}(\mathbf e_\iota/2)\),
Hermitian, \([T^1_{\rm LL},T^2_{\rm LL}]=iT^3_{\rm LL}\) and cyclically,
and the isospin Casimir
\((T_{\rm LL})^2=\sum_\iota(T^\iota_{\rm LL})^2\) & new &
\texttt{T\_LL\^{}iota} \texttt{T\_LL\^{}3} \texttt{T\_LL\^{}2} \\
\(\mathrm{Cas}_{\rm LL}\) & the Casimir operator of the rotations,
\(\mathrm{Cas}_{\rm LL}=-\sum_jD_{\rm LL}(J^j_{\rm LL})^2\); its
eigenvalue is 0 for spin 0, 3/4 for spin 1/2 and 2 for spin 1 & new &
\texttt{Cas\_LL} \\
\end{longtable}

\textbf{4 The covariant families and their classification}

\begin{longtable}[]{@{}
  >{\raggedright\arraybackslash}p{(\columnwidth - 6\tabcolsep) * \real{0.2200}}
  >{\raggedright\arraybackslash}p{(\columnwidth - 6\tabcolsep) * \real{0.5000}}
  >{\raggedright\arraybackslash}p{(\columnwidth - 6\tabcolsep) * \real{0.0800}}
  >{\raggedright\arraybackslash}p{(\columnwidth - 6\tabcolsep) * \real{0.2000}}@{}}
\toprule\noalign{}
\begin{minipage}[b]{\linewidth}\raggedright
symbol
\end{minipage} & \begin{minipage}[b]{\linewidth}\raggedright
meaning
\end{minipage} & \begin{minipage}[b]{\linewidth}\raggedright
origin
\end{minipage} & \begin{minipage}[b]{\linewidth}\raggedright
plain-text names
\end{minipage} \\
\midrule\noalign{}
\endhead
\bottomrule\noalign{}
\endlastfoot
\(\mathcal F^{\rm A}_{\rm LL}\), \(\mathcal F^{\rm B}_{\rm LL}\),
\(\mathcal F^{\rm C}_{\rm LL}\) & the three families of covariant
brackets: \(\mathcal F^{\rm A}_{\rm LL}\) on
\(\mathfrak l_{\rm LL}\oplus\mathfrak s_{\rm LL}\) (dimension 14),
\(\mathcal F^{\rm B}_{\rm LL}\) on
\(\mathfrak l_{\rm LL}\oplus\mathfrak p_{\rm LL}\oplus\mathfrak s_{\rm LL}\)
(18), \(\mathcal F^{\rm C}_{\rm LL}\) on
\(\mathfrak l_{\rm LL}\oplus\mathfrak{su}(2)\oplus\mathfrak s_{\rm LL}\)
(17) & new & \texttt{F\_A} \texttt{F\_B} \texttt{F\_C} \\
\(\mathfrak p_{\rm LL}\); \(\mathbf r\), \(\mathbf r'\) & the
translations of the family \(\mathcal F^{\rm B}_{\rm LL}\): the real
space of the Hermitian bi-quaternions \(\mathbf r=\mathbf r^\dagger\)
(dimension 4), on which \(\mathfrak l_{\rm LL}\) acts by
\([\mathbf u,\mathbf r]=\mathbf u\mathbf r+\mathbf r\mathbf u^\dagger\)
& new & \texttt{p\_LL} \texttt{rP} \texttt{rP\textquotesingle{}} \\
\(c^{\rm tw}_{\rm LL}\), \(c^{\rm sl}_{\rm LL}\),
\(c^{\rm PP}_{\rm LL}\), \(c^{\rm Ps}_{\rm LL}\),
\(c^{\rm sP}_{\rm LL}\), \(c^{\rm TT}_{\rm LL}\),
\(c^{\rm Ts}_{\rm LL}\), \(c^{\rm sT}_{\rm LL}\) & the coefficients of
the families:
\([\mathbf u,\mathbf b\mathbf e_4]=c^{\rm tw}_{\rm LL}(W_{\rm LL}(\mathbf u)\mathbf b)\mathbf e_4\)
(the twist),
\([\mathbf b\mathbf e_4,\mathbf b'\mathbf e_4]\supset c^{\rm sl}_{\rm LL}(\bar{\mathbf b}\mathbf b'-\bar{\mathbf b}'\mathbf b)\),
\([\mathbf r,\mathbf r']=c^{\rm PP}_{\rm LL}(\mathbf r\bar{\mathbf r}'-\mathbf r'\bar{\mathbf r})\),
\([\mathbf r,\mathbf b\mathbf e_4]=c^{\rm Ps}_{\rm LL}(\mathbf b^*\bar{\mathbf r})\mathbf e_4\),
\([\mathbf b\mathbf e_4,\mathbf b'\mathbf e_4]\supset c^{\rm sP}_{\rm LL}\,i(\bar{\mathbf b}\mathbf b'^*-\bar{\mathbf b}'\mathbf b^*)\),
\([\mathbf t,\mathbf t']=c^{\rm TT}_{\rm LL}(\mathbf t\mathbf t'-\mathbf t'\mathbf t)\),
\([\mathbf t,\mathbf b\mathbf e_4]=c^{\rm Ts}_{\rm LL}(\mathbf t\mathbf b)\mathbf e_4\),
\([\mathbf b\mathbf e_4,\mathbf b'\mathbf e_4]\supset\mathrm{Re}[c^{\rm sT}_{\rm LL}(\mathbf b\bar{\mathbf b}'-\mathbf b'\bar{\mathbf b})]\);
complex except \(c^{\rm sP}_{\rm LL}\), \(c^{\rm TT}_{\rm LL}\) (real);
plain text \texttt{ctw\_re}, \texttt{ctw\_im}: real and imaginary part
of \(c^{\rm tw}_{\rm LL}\), and likewise & new & \texttt{c\_tw}
\texttt{c\_sl} \texttt{c\_PP} \texttt{c\_Ps} \texttt{c\_sP}
\texttt{c\_TT} \texttt{c\_Ts} \texttt{c\_sT} \texttt{csP} \texttt{cTT}
\texttt{ctw\_re} \texttt{ctw\_im} \texttt{csl\_re} \texttt{csl\_im}
\texttt{cPP\_re} \texttt{cPP\_im} \texttt{cPs\_re} \texttt{cPs\_im}
\texttt{cTs\_re} \texttt{cTs\_im} \texttt{csT\_re} \texttt{csT\_im} \\
\(\lambda^{\rm sc}_{\rm LL}\) & a complex number \(\neq0\) by which the
spinor charges are rescaled,
\(\mathbf b\mapsto\lambda^{\rm sc}_{\rm LL}\mathbf b\) & new &
\texttt{lambda\_sc} \\
\(\mathcal I^{\rm Jac}_{\rm LL}\), \(\mathcal I^{\rm Mal}_{\rm LL}\);
\(w^{\rm sat}_{\rm LL}\) & the ideals of polynomials in the real
coefficients of a family generated by the Jacobi and by the Malcev
identity; \(w^{\rm sat}_{\rm LL}\) the auxiliary variable of the
saturation (the points of the variety at which a given polynomial does
not vanish are the points of the variety of the ideal enlarged by that
polynomial times \(w^{\rm sat}_{\rm LL}\), minus 1) & new &
\texttt{w\_sat} \\
\end{longtable}

\textbf{5 The manifolds}

\begin{longtable}[]{@{}
  >{\raggedright\arraybackslash}p{(\columnwidth - 6\tabcolsep) * \real{0.2200}}
  >{\raggedright\arraybackslash}p{(\columnwidth - 6\tabcolsep) * \real{0.5000}}
  >{\raggedright\arraybackslash}p{(\columnwidth - 6\tabcolsep) * \real{0.0800}}
  >{\raggedright\arraybackslash}p{(\columnwidth - 6\tabcolsep) * \real{0.2000}}@{}}
\toprule\noalign{}
\begin{minipage}[b]{\linewidth}\raggedright
symbol
\end{minipage} & \begin{minipage}[b]{\linewidth}\raggedright
meaning
\end{minipage} & \begin{minipage}[b]{\linewidth}\raggedright
origin
\end{minipage} & \begin{minipage}[b]{\linewidth}\raggedright
plain-text names
\end{minipage} \\
\midrule\noalign{}
\endhead
\bottomrule\noalign{}
\endlastfoot
\(\mathbb R^{1,7}_{\rm LL}\), \(\mathbb R^{1,3}_{\rm LL}\);
\(x^p_{\rm LL}\) & the Hermitian bi-octonions
\(\mathbf x=\mathbf x^\dagger=x^0_{\rm LL}\mathbf e_0+i\sum_{p=1}^7x^p_{\rm LL}\mathbf e_p\)
(real \(x^p_{\rm LL}\)): the \(7+1\) Minkowski space with the real norm
\(N_{\rm LL}(\mathbf x)=(x^0_{\rm LL})^2-\sum_p(x^p_{\rm LL})^2\); and
its subspace of Hermitian bi-quaternions (\(p\le3\)), Minkowski space of
\(3+1\) dimensions & new & \texttt{R\^{}(1,7)\_LL}
\texttt{R\^{}(1,3)\_LL} \texttt{x\_LL\^{}0} \texttt{x\_LL\^{}p}
\texttt{x\_LL\^{}1} \texttt{x\_LL\^{}2} \texttt{x\_LL\^{}3}
\texttt{x\_LL\^{}j} \\
\(\mathbb O^{\rm s}_{\rm LL}\) & the split octonions inside
\(\mathbb B_{\rm LL}\): the real space
\(\mathbb H+(i\mathbb H)\mathbf e_4\) (real combinations of
\(\mathbf e_0,\dots,\mathbf e_3,i\mathbf e_4,\dots,i\mathbf e_7\)),
closed under the product, on which \(N_{\rm LL}\) has signature
\((4,4)\): the \(4+4\) manifold (plain text \texttt{R\^{}(4,4)}) & new &
\texttt{Os} \texttt{R\^{}(4,4)} \\
\(S^7_{\rm LL}\), \(S^{3,4}_{\rm LL}\) & the unit elements
(\(N_{\rm LL}=1\)) of \(\mathbb O\) (the seven-sphere) and of
\(\mathbb O^{\rm s}_{\rm LL}\) (the pseudo-sphere): the two real forms
of \(\mathbb L_{\rm LL}\), both Moufang loops & new & \texttt{S\^{}7}
\texttt{S\^{}(3,4)} \\
\(X^{\rm LL}_{\mathbf x}\) & the frame field
\(X^{\rm LL}_{\mathbf x}(\mathbf w)=\mathbf w\mathbf x\) on
\(\mathbb L_{\rm LL}\) for imaginary \(\mathbf x\) (a vector field: the
velocity of the curve \(\mathbf w(\theta_{\rm LL})\) with
\(d\mathbf w/d\theta_{\rm LL}=\mathbf w\mathbf x\)) & new &
\texttt{X\_x} \texttt{X\_y} \\
\(\mathsf L_{\rm LL}(\mathbf g)\) & the real \(4\times4\) matrix of the
Lorentz transformation
\(\mathbf x\mapsto\mathbf g\mathbf x\mathbf g^\dagger\) of
\(\mathbb R^{1,3}_{\rm LL}\) in the coordinates
\(x^0_{\rm LL},\dots,x^3_{\rm LL}\), for a unit bi-quaternion
\(\mathbf g\): its entry in row \(p\) and column \(p'\)
(\(p,p'=0,\dots,3\)) is
\(\tfrac12\,\mathrm{Re}\,\mathrm{tr}(\sigma^{\rm LL}_pM_{\rm LL}(\mathbf g)\,\sigma^{\rm LL}_{p'}M_{\rm LL}(\mathbf g)^\dagger)\),
with \(\sigma^{\rm LL}_0\) the unit matrix & new & \texttt{Lmat} \\
\(\chi_{\rm LL}\) & a two-component spinor (a column of two complex
numbers) on which a unit bi-quaternion \(\mathbf g\) acts by
\(\chi_{\rm LL}\mapsto M_{\rm LL}(\mathbf g)\chi_{\rm LL}\); the
Hermitian matrix \(\chi_{\rm LL}\chi_{\rm LL}^\dagger\) is the matrix
\(M_{\rm LL}(\mathbf x)\) of a point \(\mathbf x\) of the future light
cone of \(\mathbb R^{1,3}_{\rm LL}\) & new & \texttt{chi} \\
\(\theta_{\rm LL}\) & the real parameter of a one-parameter family of
transformations or of a curve (an angle, a rapidity) & new &
\texttt{theta} \\
\(\mathcal V_{\rm LL}(\mathbf u)\) & the vector action of a Lorentz
element on \(\mathbb R^{1,3}_{\rm LL}\),
\(\mathcal V_{\rm LL}(\mathbf u)\mathbf x=\mathbf u\mathbf x+\mathbf x\mathbf u^\dagger\),
a real \(4\times4\) matrix with entries
\(\mathcal V_{\rm LL}(\mathbf u)^p{}_{p'}\) in the coordinates
\(x^p_{\rm LL}\) (plain text \texttt{V\_LL\^{}p\_p\textquotesingle{}}) &
new & \texttt{V\_LL} \\
\end{longtable}

\textbf{6 The spin-zero superfluid}

\begin{longtable}[]{@{}
  >{\raggedright\arraybackslash}p{(\columnwidth - 6\tabcolsep) * \real{0.2200}}
  >{\raggedright\arraybackslash}p{(\columnwidth - 6\tabcolsep) * \real{0.5000}}
  >{\raggedright\arraybackslash}p{(\columnwidth - 6\tabcolsep) * \real{0.0800}}
  >{\raggedright\arraybackslash}p{(\columnwidth - 6\tabcolsep) * \real{0.2000}}@{}}
\toprule\noalign{}
\begin{minipage}[b]{\linewidth}\raggedright
symbol
\end{minipage} & \begin{minipage}[b]{\linewidth}\raggedright
meaning
\end{minipage} & \begin{minipage}[b]{\linewidth}\raggedright
origin
\end{minipage} & \begin{minipage}[b]{\linewidth}\raggedright
plain-text names
\end{minipage} \\
\midrule\noalign{}
\endhead
\bottomrule\noalign{}
\endlastfoot
\(M_{\rm LL}(\mathbf b)\) & the complex \(2\times2\) matrix of a
bi-quaternion, \(M_{\rm LL}(\mathbf e_0)=\sigma^{\rm LL}_0\),
\(M_{\rm LL}(\mathbf e_j)=-i\sigma^{\rm LL}_j\), extended
complex-linearly;
\(M_{\rm LL}(\mathbf b\mathbf b')=M_{\rm LL}(\mathbf b)M_{\rm LL}(\mathbf b')\),
\(\det M_{\rm LL}(\mathbf b)=N_{\rm LL}(\mathbf b)\) (plain text
\texttt{M\_2(C)}: all complex \(2\times2\) matrices) & new &
\texttt{M\_LL} \texttt{M\_2(C)} \\
\(\psi'_{\rm LL}\) & the Weyl isodoublet of a spinor charge
\(\mathbf b\mathbf e_4\): the \(2\times2\) matrix
\(\psi'_{\rm LL}=\mathsf E_{\rm LL}^{-1}M_{\rm LL}(\mathbf b)^T\), whose
first column (isospin up) and second column (isospin down) are Weyl
spinors & new & \texttt{psi\textquotesingle{}\_LL} \\
\(\Psi_{\rm LL}\) & the Dirac isodoublet of the spinor charges: a field
of 8 complex Grassmann components (Dirac \(\otimes\) isospin), whose
\(\gamma^{\rm ch}_{\rm LL}=+1\) half is built from
\(\mathfrak s_{\rm LL}\) and whose other half from the complex-conjugate
super-loop & new & \texttt{Psi\_LL} \texttt{Psi\_LL\^{}T}
\texttt{Psi\_LL\^{}dag} \\
\(\gamma^p_{\rm LL}\), \(\gamma^{\rm ch}_{\rm LL}\), \(C_{\rm LL}\),
\(\alpha^j_{\rm LL}\), \(\beta_{\rm LL}\) & the \(4\times4\) Dirac
matrices of \(\mathbb R^{1,3}_{\rm LL}\) in the chiral basis,
\(\gamma^0_{\rm LL}=\begin{pmatrix}0&\sigma^{\rm LL}_0\\\sigma^{\rm LL}_0&0\end{pmatrix}\),
\(\gamma^j_{\rm LL}=\begin{pmatrix}0&\sigma^{\rm LL}_j\\-\sigma^{\rm LL}_j&0\end{pmatrix}\);
the chirality matrix
\(\gamma^{\rm ch}_{\rm LL}=i\gamma^0_{\rm LL}\gamma^1_{\rm LL}\gamma^2_{\rm LL}\gamma^3_{\rm LL}=\mathrm{diag}(-1,-1,1,1)\);
the charge-conjugation matrix
\(C_{\rm LL}=i\gamma^2_{\rm LL}\gamma^0_{\rm LL}\);
\(\alpha^j_{\rm LL}=\gamma^0_{\rm LL}\gamma^j_{\rm LL}\),
\(\beta_{\rm LL}=\gamma^0_{\rm LL}\) (on Dirac \(\otimes\) isospin each
acts as \(\otimes\sigma^{\rm LL}_0\)) & new & \texttt{gamma\_LL\^{}p}
\texttt{gamma\_LL\^{}p\textquotesingle{}} \texttt{gamma\_ch}
\texttt{C\_LL} \texttt{alpha\_LL\^{}j} \texttt{beta\_LL} \\
\(\mathsf K^\iota_{\rm LL}\), \(\mathsf K'^\iota_{\rm LL}\),
\(\mathsf K^{\rm s}_{\rm LL}\) & the \(8\times8\) pair matrices of the
spin-zero channels: scalar
\(\mathsf K^\iota_{\rm LL}=C_{\rm LL}\gamma^{\rm ch}_{\rm LL}\otimes\mathsf E_{\rm LL}\tau^{\rm LL}_\iota\)
(\(0^+\)), pseudoscalar
\(\mathsf K'^\iota_{\rm LL}=C_{\rm LL}\otimes\mathsf E_{\rm LL}\tau^{\rm LL}_\iota\)
(\(0^-\)), and the \(16\times16\) singlet of two species
\(\mathsf K^{\rm s}_{\rm LL}=C_{\rm LL}\gamma^{\rm ch}_{\rm LL}\otimes\mathsf E_{\rm LL}\otimes\mathsf E_{\rm LL}\)
(the last factor on the species index); the plain-text name Kch alone,
and the letter \(\mathsf K\) without scripts (in Theorem SF1 and in the
explanation of the folder explanations), stand for any pair matrix
(\(8\times8\) for one field, \(16\times16\) for two species), the pair
of two spinor charges being \(\Psi^T_{\rm LL}\mathsf K\Psi_{\rm LL}\);
\(\mathsf K_{12}\) is the entry of \(\mathsf K\) in row 1 and column 2 &
new & \texttt{Kch} \texttt{Kch\_iota} \texttt{Kch\_1} \texttt{Kch\_2}
\texttt{Kch\_3} \texttt{Kch\textquotesingle{}\_iota}
\texttt{Kch\_singlet} \\
\(\Pi^\iota_{\rm LL}\) & the pair operators
\(\Pi^\iota_{\rm LL}=\sum_{\vec q_{\rm LL}}\Psi_{\rm LL}(\vec q_{\rm LL})^T\mathsf K^\iota_{\rm LL}\Psi_{\rm LL}(-\vec q_{\rm LL})\),
which annihilate two spinor charges of opposite momenta & new &
\texttt{Pi\_iota} \\
\(\vec q_{\rm LL}\), \(q_{\rm LL}\), \(q^j_{\rm LL}\) & the momentum of
a particle in the space of \(\mathbb R^{1,3}_{\rm LL}\), its length and
its components & new & \texttt{q} \\
\(m_{\rm LL}\), \(\mu_{\rm LL}\), \(\Lambda_{\rm LL}\) & the mass of the
spinor charges, the chemical potential, and the momentum cutoff
(\(q_{\rm LL}<\Lambda_{\rm LL}\); \(\Lambda_{\rm LL}=1\) sets the units,
except where it varies) & new & \texttt{m} \texttt{mu}
\texttt{Lambda} \\
\(\varepsilon_{\rm LL}\), \(E^\pm_{\rm LL}\), \(\tilde E^\pm_{\rm LL}\)
& the energy \(\varepsilon_{\rm LL}=(q_{\rm LL}^2+m_{\rm LL}^2)^{1/2}\),
the quasi-particle energies of the scalar channel
\(E^\pm_{\rm LL}=[(\varepsilon_{\rm LL}\pm\mu_{\rm LL})^2+\Delta_{\rm LL}^2]^{1/2}\),
a function of the gap (written \(E^\pm_{\rm LL}(\Delta^+_{\rm LL})\) and
\(E^\pm_{\rm LL}(\Delta^-_{\rm LL})\) when \(\Delta_{\rm LL}\) is
replaced by one of the two gaps of a general state), and those of the
pseudoscalar channel,
\((\tilde E^\pm_{\rm LL})^2=\varepsilon_{\rm LL}^2+\mu_{\rm LL}^2+\Delta_{\rm LL}^2\pm2[\mu_{\rm LL}^2\varepsilon_{\rm LL}^2+m_{\rm LL}^2\Delta_{\rm LL}^2]^{1/2}\)
& new & \texttt{eps} \texttt{E\^{}+} \texttt{E\^{}-} \texttt{E\^{}(+-)}
\texttt{E\_minus\_polar} \texttt{E\_plus\_polar} \texttt{Etilde\^{}+}
\texttt{Etilde\^{}-} \texttt{Etilde\^{}(+-)} \texttt{Etilde\_minus}
\texttt{Etilde\_plus} \\
\(G_{\rm LL}\), \(G^{\rm crit}_{\rm LL}\) & the coupling of the
attractive spin-zero contact interaction (dimension mass\(^{-2}\)) and
its critical value at \(\mu_{\rm LL}=0\), above which the vacuum
condenses & new & \texttt{G} \texttt{G\_crit}
\texttt{G\_over\_Gcrit\_m0} \texttt{G\_crit\_mu0\_m0} \\
\(\mathcal H^{\rm sf}_{\rm LL}\) & the Hamiltonian of the superfluid,
\(\sum_{\vec q_{\rm LL}}\Psi^\dagger_{\rm LL}(\mathsf h_{\rm LL}-\mu_{\rm LL})\Psi_{\rm LL}-\tfrac14G_{\rm LL}\sum_\iota(\Pi^\iota_{\rm LL})^\dagger\Pi^\iota_{\rm LL}\)
& new & \texttt{Hsf} \\
\(\mathsf h_{\rm LL}(\vec q_{\rm LL})\) & the \(8\times8\) Hamiltonian
matrix of a plane wave,
\(\sum_j\alpha^j_{\rm LL}q^j_{\rm LL}+\beta_{\rm LL}m_{\rm LL}\) (times
\(\sigma^{\rm LL}_0\) on the isospin index); eigenvalues
\(\pm\varepsilon_{\rm LL}\) & new & \texttt{h(q)} \\
\(\vec d_{\rm LL}\), \(d^\iota_{\rm LL}\); \(\zeta_{\rm LL}\) & the
order parameter, the isospin vector
\(\vec d_{\rm LL}=(d^1_{\rm LL},d^2_{\rm LL},d^3_{\rm LL})\in\mathbb C^3\)
with
\(d^\iota_{\rm LL}=-\tfrac12G_{\rm LL}\langle\Pi^\iota_{\rm LL}\rangle\),
and its non-unitarity
\(\zeta_{\rm LL}=\lvert\vec d_{\rm LL}\times\vec d^{\,*}_{\rm LL}\rvert/\lvert\vec d_{\rm LL}\rvert^2\in[0,1]\)
(plain text \texttt{cross(dvec,\ dvec\^{}*)}: the cross product) & new &
\texttt{dvec} \texttt{d\_iota} \texttt{abs\_dvec\_at\_min}
\texttt{zeta} \\
\(\Delta_{\rm LL}\), \(\Delta^\pm_{\rm LL}\) & the gap:
\(\Delta_{\rm LL}=\lvert\vec d_{\rm LL}\rvert\) for a unitary state, and
the two gaps
\((\Delta^\pm_{\rm LL})^2=\lvert\vec d_{\rm LL}\rvert^2(1\pm\zeta_{\rm LL})\)
of a general state & new & \texttt{Delta} \texttt{Delta\_+}
\texttt{Delta\_-} \texttt{Delta\_(+-)} \\
\(\mathsf P_{\rm LL}\); \(\mathsf H_{\rm LL}(\vec q_{\rm LL})\);
\(\mathcal N_{\rm LL}(\vec q_{\rm LL})\) & the pairing block
\(\mathsf P_{\rm LL}=\sum_\iota d^\iota_{\rm LL}(\mathsf K^\iota_{\rm LL})^\dagger\),
the \(16\times16\) Bogoliubov-de Gennes matrix
\(\mathsf H_{\rm LL}=\begin{pmatrix}\mathsf h_{\rm LL}(\vec q_{\rm LL})-\mu_{\rm LL}&\mathsf P_{\rm LL}\\\mathsf P_{\rm LL}^\dagger&-(\mathsf h_{\rm LL}(-\vec q_{\rm LL})-\mu_{\rm LL})^T\end{pmatrix}\)
and the Nambu spinor
\(\mathcal N_{\rm LL}=(\Psi_{\rm LL}(\vec q_{\rm LL}),\Psi_{\rm LL}^{*}(-\vec q_{\rm LL}))\)
on which it acts & new & \texttt{Ppair} \texttt{H\_BdG} \\
\(\mathsf D_{\rm LL}\) & the isospin part of the pairing block of the
scalar channel, the \(2\times2\) matrix
\(\mathsf D_{\rm LL}=\sum_\iota d^\iota_{\rm LL}(\mathsf E_{\rm LL}\tau^{\rm LL}_\iota)^\dagger\),
so that
\(\mathsf P_{\rm LL}=(C_{\rm LL}\gamma^{\rm ch}_{\rm LL})^\dagger\otimes\mathsf D_{\rm LL}\);
it is also the isospin wavefunction of the condensed pairs (rows and
columns: isospin up, down); plain text \texttt{Dpair} (not
\texttt{D\_LL}, which is the Lorentz derivation) & new &
\texttt{Dpair} \\
\(\Omega_{\rm LL}\) & the grand potential per volume at zero
temperature, measured from the normal state:
\(\Omega_{\rm LL}=\lvert\vec d_{\rm LL}\rvert^2/G_{\rm LL}-\tfrac12\int\!d^3q_{\rm LL}(2\pi)^{-3}\sum(\text{positive eigenvalues of }\mathsf H_{\rm LL})\)
minus its value at \(\vec d_{\rm LL}=0\) & new & \texttt{Omega}
\texttt{Omega\_min} \texttt{Omega\_minus\_Omega0}
\texttt{Omega\_at\_polar\_gap} \\
\(k^{\rm F}_{\rm LL}\), \(\nu^{\rm F}_{\rm LL}\) & the Fermi momentum
\((\mu_{\rm LL}^2-m_{\rm LL}^2)^{1/2}\) and the density of states of one
species and one spin state at the Fermi surface,
\(\nu^{\rm F}_{\rm LL}=k^{\rm F}_{\rm LL}\mu_{\rm LL}/(2\pi^2)\) & new &
\texttt{kF} \texttt{nuF} \\
\(F_{\rm LL}\); \(U(1)^{\rm F}_{\rm LL}\); \(SU(2)_{\rm v}\) & the
generator of the fermion number
\(\Psi_{\rm LL}\to e^{i\theta_{\rm LL}}\Psi_{\rm LL}\) (it multiplies
\(\vec d_{\rm LL}\) by \(2i\)) and its group; the isospin group, which
rotates the isospin index of \(\Psi_{\rm LL}\) and \(\vec d_{\rm LL}\)
as a vector & new & \texttt{F\_LL} \texttt{U(1)\_F\_LL}
\texttt{SU(2)\_v} \\
\(\Psi^{(1)}_{\rm LL}\), \(\Psi^{(2)}_{\rm LL}\) & two species of the
field \(\Psi_{\rm LL}\), combined into one field of 16 components with a
species index & new & \\
\end{longtable}

\textbf{7 Power counting and renormalizability}

\begin{longtable}[]{@{}
  >{\raggedright\arraybackslash}p{(\columnwidth - 6\tabcolsep) * \real{0.2200}}
  >{\raggedright\arraybackslash}p{(\columnwidth - 6\tabcolsep) * \real{0.5000}}
  >{\raggedright\arraybackslash}p{(\columnwidth - 6\tabcolsep) * \real{0.0800}}
  >{\raggedright\arraybackslash}p{(\columnwidth - 6\tabcolsep) * \real{0.2000}}@{}}
\toprule\noalign{}
\begin{minipage}[b]{\linewidth}\raggedright
symbol
\end{minipage} & \begin{minipage}[b]{\linewidth}\raggedright
meaning
\end{minipage} & \begin{minipage}[b]{\linewidth}\raggedright
origin
\end{minipage} & \begin{minipage}[b]{\linewidth}\raggedright
plain-text names
\end{minipage} \\
\midrule\noalign{}
\endhead
\bottomrule\noalign{}
\endlastfoot
\([\ \cdot\ ]\) (one argument) & the mass dimension of a field or a
coupling, in units in which \(\hbar\) and the speed of light are 1 (a
bracket with two or three arguments is a commutator or an associator) &
new & \texttt{{[}Psi\_LL{]}} \texttt{{[}phi{]}} \texttt{{[}G{]}}
\texttt{{[}g\_Y{]}} \texttt{{[}lam4{]}} \\
\(\phi_{\rm LL}\), \(g^{\rm Y}_{\rm LL}\), \(\lambda^{(1)}_{\rm LL}\),
\(\lambda^{(2)}_{\rm LL}\) & the Yukawa completion: a complex
isospin-triplet pair field
\(\phi_{\rm LL}=(\phi^1_{\rm LL},\phi^2_{\rm LL},\phi^3_{\rm LL})\)
coupled to \(\Psi^T_{\rm LL}\mathsf K^\iota_{\rm LL}\Psi_{\rm LL}\) with
the Yukawa coupling \(g^{\rm Y}_{\rm LL}\), and the couplings
\(\lambda^{(1)}_{\rm LL}\), \(\lambda^{(2)}_{\rm LL}\) of the quartic
terms \((\phi^\dagger_{\rm LL}\phi_{\rm LL})^2\) and
\(\lvert\phi^T_{\rm LL}\phi_{\rm LL}\rvert^2\) (plain text
\texttt{lam4}: either of them) & new & \texttt{phi} \texttt{g\_Y}
\texttt{lam4} \texttt{lam4\_1} \texttt{lam4\_2} \\
\(\mathcal L^{\rm Y}_{\rm LL}\), \(\mathcal L^{\rm c}_{\rm LL}\) & the
Lagrangian densities of the Yukawa form and of the contact form of the
field theory of the spinor charges on \(\mathbb R^{1,3}_{\rm LL}\) & new
& \texttt{Lagr\_Y} \texttt{Lagr\_c} \\
\(m^\phi_{\rm LL}\) & the mass of the pair field \(\phi_{\rm LL}\) in
the Yukawa form & new & \texttt{m\_phi} \\
\(I^{\rm ren}_{\rm LL}\) & the renormalized one-loop bubble,
\(I^{\rm ren}_{\rm LL}(p^{\rm E}_{\rm LL})=I^{\rm bub}_{\rm LL}(p^{\rm E}_{\rm LL})-I^{\rm bub}_{\rm LL}(0)-(p^{\rm E}_{\rm LL})^2\,\partial I^{\rm bub}_{\rm LL}/\partial(p^{\rm E}_{\rm LL})^2\)
(the derivative taken at \(p^{\rm E}_{\rm LL}=0\)): the bubble minus its
mass and wave-function counterterms & new & \texttt{Iren} \\
\(\deg_{\rm LL}\) & the superficial degree of divergence of a
one-particle-irreducible graph: the power of the cutoff with which it
grows when all its loop momenta grow together & new & \texttt{deg} \\
\(n^{\rm loop}_{\rm LL}\), \(n^{\rm int,F}_{\rm LL}\),
\(n^{\rm int,B}_{\rm LL}\), \(n^{\rm ext,F}_{\rm LL}\),
\(n^{\rm ext,B}_{\rm LL}\), \(n^{\rm v}_{\rm LL}\),
\(n^{\rm v,Y}_{\rm LL}\), \(n^{\rm v,4}_{\rm LL}\) & the numbers of
loops, of internal fermion and boson lines, of external fermion and
boson lines, of vertices, of Yukawa vertices and of quartic vertices of
a graph & new & \texttt{nloop} \texttt{nint\_F} \texttt{nint\_B}
\texttt{next\_F} \texttt{next\_B} \texttt{nvert} \texttt{nvert\_Y}
\texttt{nvert\_4} \\
\(I^{\rm bub}_{\rm LL}\); \(k^{\rm E}_{\rm LL}\), \(p^{\rm E}_{\rm LL}\)
& the one-loop bubble of the Yukawa completion,
\(I^{\rm bub}_{\rm LL}(p^{\rm E}_{\rm LL})=\int_{k^{\rm E}_{\rm LL}<\Lambda_{\rm LL}}\frac{d^4k^{\rm E}_{\rm LL}}{(2\pi)^4}\frac{k^{\rm E}_{\rm LL}\cdot(k^{\rm E}_{\rm LL}+p^{\rm E}_{\rm LL})+m_{\rm LL}^2}{[(k^{\rm E}_{\rm LL})^2+m_{\rm LL}^2][(k^{\rm E}_{\rm LL}+p^{\rm E}_{\rm LL})^2+m_{\rm LL}^2]}\),
with the Euclidean loop and external four-momenta
\(k^{\rm E}_{\rm LL}\), \(p^{\rm E}_{\rm LL}\) (and their lengths) & new
& \texttt{Ibub} \texttt{Ibub\_0} \texttt{Ibub\_pE}
\texttt{Ibub\_pE\_minus\_Ibub\_0} \texttt{kE} \texttt{pE} \\
\end{longtable}

\textbf{8 Groups, Lie algebras, operations and labels}

\begin{longtable}[]{@{}
  >{\raggedright\arraybackslash}p{(\columnwidth - 6\tabcolsep) * \real{0.2200}}
  >{\raggedright\arraybackslash}p{(\columnwidth - 6\tabcolsep) * \real{0.5000}}
  >{\raggedright\arraybackslash}p{(\columnwidth - 6\tabcolsep) * \real{0.0800}}
  >{\raggedright\arraybackslash}p{(\columnwidth - 6\tabcolsep) * \real{0.2000}}@{}}
\toprule\noalign{}
\begin{minipage}[b]{\linewidth}\raggedright
symbol
\end{minipage} & \begin{minipage}[b]{\linewidth}\raggedright
meaning
\end{minipage} & \begin{minipage}[b]{\linewidth}\raggedright
origin
\end{minipage} & \begin{minipage}[b]{\linewidth}\raggedright
plain-text names
\end{minipage} \\
\midrule\noalign{}
\endhead
\bottomrule\noalign{}
\endlastfoot
\(SL(2,\mathbb C)\), \(Spin(3,1)^+\), \(SO(4,\mathbb C)\),
\(\mathrm G_2\), \(\mathrm G_2(\mathbb C)\), \(SO(1,7)\), \(Spin(7)\),
\(Spin(7,\mathbb C)\), \(SL(2,\mathbb R)\), \(Spin(2,1)\), \(SU(2)\),
\(U(1)\), \(\mathbb Z_2\) & groups: the complex \(2\times2\) matrices of
determinant 1 (the unit bi-quaternions), equal to the connected double
cover of the Lorentz group of \(3+1\) dimensions; the complex orthogonal
group of four dimensions; the compact exceptional group \(\mathrm G_2\)
(automorphisms of \(\mathbb O\)) and its complex form (automorphisms of
\(\mathbb B_{\rm LL}\)); the Lorentz group of \(7+1\) dimensions; the
spin groups of seven dimensions (compact and complex); the real
\(2\times2\) matrices of determinant 1, the double cover of the Lorentz
group of \(2+1\) dimensions; the special unitary group of two
dimensions; the phases; the group of two elements & new &
\texttt{SL(2,C)} \texttt{Spin(3,1)\^{}+} \texttt{SO(4,C)} \texttt{G2}
\texttt{G2(C)} \texttt{SO(1,7)} \texttt{SO(3,1)} \texttt{Spin(7)}
\texttt{Spin(7,C)} \texttt{SL(2,R)} \texttt{Spin(2,1)} \texttt{SU(2)}
\texttt{U(1)} \texttt{Z2} \\
\(\mathfrak{so}(3,1)\), \(\mathfrak{so}(1,7)\), \(\mathfrak{so}(4,1)\),
\(\mathfrak{so}(8)\), \(\mathfrak{so}(4,4)\),
\(\mathfrak{so}(8,\mathbb C)\), \(\mathfrak{so}(3)\),
\(\mathfrak{so}(4)\), \(\mathfrak{so}(2,1)\),
\(\mathfrak{sl}(2,\mathbb R)\), \(\mathfrak{sl}(2)\),
\(\mathfrak{su}(2)\), \(\mathfrak{spin}(7)\), \(\mathfrak g_2\),
\(\mathfrak g_2(\mathbb C)\) & Lie algebras: of the Lorentz groups of
\(3+1\) and \(7+1\) dimensions; of the de Sitter group (signature of the
Killing form: 4 positive, 6 negative); of the linear maps that preserve
a definite real, a \((4,4)\) and a complex quadratic form in eight
dimensions; of the rotations of three-dimensional space and of
four-dimensional Euclidean space; of the Lorentz group of \(2+1\)
dimensions, equal to \(\mathfrak{sl}(2,\mathbb R)\);
\(\mathfrak{sl}(2)\) the simple Lie algebra of dimension 3 over any
field; \(\mathfrak{su}(2)\) isospin; \(\mathfrak{spin}(7)\) of dimension
21; the exceptional Lie algebra \(\mathfrak g_2\) (compact, dimension
14) and its complex form (real dimension 28) & new & \texttt{so(3,1)}
\texttt{so(1,7)} \texttt{so(4,1)} \texttt{so(8)} \texttt{so(4,4)}
\texttt{so(8,C)} \texttt{so(3)} \texttt{so(4)} \texttt{so(2,1)}
\texttt{sl(2,R)} \texttt{sl(2)} \texttt{su(2)} \texttt{spin(7)}
\texttt{g2} \texttt{g2(C)} \\
\(\mathrm{tr}\), \(\det\), \(\mathrm{diag}\), \(\mathrm{Re}\),
\(\mathrm{Im}\) & trace, determinant, the diagonal matrix with the
listed entries, real part, imaginary part (also: the imaginary part of
an algebra, the span of its imaginary units); plain text \texttt{exp},
\texttt{ln}, \texttt{sqrt}, \texttt{sum}, \texttt{int}, \texttt{span},
\texttt{dim}: the exponential, natural logarithm, square root, a sum, an
integral, the set of real linear combinations, a dimension & new &
\texttt{tr} \texttt{det} \texttt{diag} \texttt{Re} \texttt{Im}
\texttt{exp} \texttt{ln} \texttt{sqrt} \texttt{sum} \texttt{int}
\texttt{span} \texttt{dim} \\
\({}^T\); \(\otimes\), \(\oplus\); \(\circ\); \(\langle\ \rangle\); h.c.
& transpose; tensor product and direct sum (plain text \texttt{(x)} for
\(\otimes\)); composition of maps; expectation value in the ground
state; the Hermitian conjugate of the preceding terms & new &
\texttt{(x)} \texttt{\^{}T} \texttt{o} \\
A-1, \ldots, R-6; G-1, \ldots, G-12; LL-S1, \ldots, LL-H1; X-1, \ldots,
X-10; Z-1, \ldots, Z-8; Theorems R, L, M, SF1, \ldots, SF4 & the
identifiers of the machine checks (letters: A associativity, L Lorentz
super-loop, Q operators and eigenvalues, C the inversion element, M
Malcev classification, G Groebner bases, N manifolds, S superfluid, R
renormalizability, LL-S symbols, LL-H self-containedness, X the examples
of the explanation of the spin-zero superfluid in the folder
explanations, Z the 3+1 manifold and the renormalizability of the Yukawa
form in the folder explanations) and the labels of the theorems proved
in the manuscript & new & \\
H-grp, H-op, H-Lie, H-gr, H-dyn & the labels of the hypotheses of the
no-go theorems: the symmetries form a group (H-grp), act as linear
operators on the space of states (H-op), have generators that form a Lie
algebra (H-Lie) or a Lie superalgebra (H-gr), and the dynamical
hypotheses on the scattering matrix (H-dyn) & new & \\
\end{longtable}

\end{document}
